\chapter{Sleepwalking}

\section*{Definition}
Sleepwalking (somnambulism) is a parasomnia involving complex, semi-purposeful motor behaviors during NREM sleep (typically stage N3 slow-wave sleep), ranging from sitting up and looking around to walking, leaving the house, or performing routine activities. The individual is difficult to arouse and has no recall of the event.

\section*{What Happens in Your Body}
Sleepwalking results from an incomplete arousal from slow-wave sleep, with the motor cortex and subcortical motor systems activated while higher cortical areas (prefrontal cortex, hippocampus) remain asleep. This produces a state dissociation between conscious awareness and motor behavior.

\section*{Causes}
\begin{itemize}
  \item \textbf{Primary:} unknown cause -- most common in children and often resolves spontaneously
  \item \textbf{Sleep deprivation:} The single strongest trigger; also sleep schedule disruption, jet lag
  \item \textbf{Genetic:} Strong familial pattern; 10--20\% in first-degree relatives if a parent was affected; 45--60\% if both parents
  \item \textbf{Medical:} Fever, illness, obstructive sleep apnea, restless legs syndrome, gastroesophageal reflux, pregnancy
  \item \textbf{Psychiatric:} PTSD, anxiety disorders, depression, stress
  \item \textbf{Substance/medication:} Alcohol (especially intoxicated sleep), sedative-hypnotics (zolpidem, zopiclone), lithium, anticholinergics, bupropion
\end{itemize}

\section*{When to See a Doctor}
See your doctor if:
\begin{itemize}
  \item Sleepwalking leads to dangerous behavior (leaving the house, driving, using sharp objects)
  \item Episodes are frequent (persisting into adulthood or several times weekly)
  \item First-time onset in older adulthood raises concern for REM behavior disorder or seizures
\end{itemize}

\section*{Related Symptoms}
\begin{itemize}
  \item Sleep terrors, confusional arousals, sleep talking, bruxism, enuresis, daytime sleepiness, REM sleep behavior disorder
\end{itemize}
\textbf{Important:} Waking a sleepwalker is not dangerous (it does not cause a heart attack) but may lead to confusion or agitation; gentle guidance back to bed is preferred.

\section*{Self-Care / Home Management}
\begin{itemize}
  \item Ensure a safe sleep environment: lock doors and windows, remove sharp objects, block stairs, hide car keys
  \item Maintain consistent sleep schedules and prioritize adequate sleep duration
  \item Avoid alcohol and heavy meals before bed
  \item Practice relaxation techniques before sleep to reduce stress
\end{itemize}

\section*{How Your Doctor Will Evaluate You}
\begin{itemize}
  \item Clinical history from patient and bed partner/family; focus on timing, frequency, behaviors, and triggers
  \item Differentiate from REM sleep behavior disorder (RBD): RBD occurs in REM sleep, involves dream enactment, and the individual is easily awakened with dream recall
  \item Video polysomnography if atypical features (onset in adulthood, bizarre or stereotyped behaviors, potential seizure)
  \item Consider EEG to exclude nocturnal frontal lobe epilepsy
\end{itemize}

\section*{Treatment Options}
\begin{itemize}
  \item Safety precautions are the mainstay of management
  \item Address triggers: treat sleep apnea, improve sleep hygiene, eliminate precipitating medications/alcohol
  \item Scheduled anticipatory awakenings: waking the child 15--30 minutes before the typical episode time
  \item Pharmacotherapy for severe or dangerous cases: low-dose clonazepam, diazepam, or melatonin
  \item Hypnosis has evidence for use in children with sleepwalking
\end{itemize}

\section*{References}
\begin{thebibliography}{99}
\bibitem{sleepwalking:aasm-sleepwalking} American Academy of Sleep Medicine. \emph{International Classification of Sleep Disorders}. 4th ed. AASM; 2023. Chapter: Sleepwalking.
\bibitem{sleepwalking:uptodate-sleepwalking} Kryger MH, et al. Sleepwalking in adults. In: UpToDate, Post TW (Ed), UpToDate, Waltham, MA; 2025.
\bibitem{sleepwalking:lancet-neurol-parasomnia} Howell MJ, et al. "Parasomnias: diagnosis and management." \emph{Lancet Neurol}. 2024;23(2):178-189.
\bibitem{sleepwalking:bmj-sleepwalking} Pressman MR, et al. "Sleepwalking." \emph{BMJ}. 2023;383:e074543.
\bibitem{sleepwalking:sleep-parasomnia} Zadra A, et al. "Disorders of arousal: clinical practice update." \emph{Sleep}. 2024;47(2):zsad312.
\end{thebibliography}
